\documentclass[10pt,a4paper]{amsart}
\usepackage{amsmath,amssymb,amsthm,graphicx}
\usepackage{verbatim}
\newenvironment{reference}{\comment}{\endcomment}
%\newenvironment{reference}{}{}
\newenvironment{slogan}{\comment}{\endcomment}
\newenvironment{history}{\comment}{\endcomment}

% For commutative diagrams we use Xy-pic
\usepackage[all]{xy}

% We use 2cell for 2-commutative diagrams.
\xyoption{2cell}
% Theorem environments.
%
\theoremstyle{plain}
\newtheorem{theorem}[subsection]{Theorem}
\newtheorem{proposition}[subsection]{Proposition}
\newtheorem{lemma}[subsection]{Lemma}

\theoremstyle{definition}
\newtheorem{definition}[subsection]{Definition}
\newtheorem{example}[subsection]{Example}
\newtheorem{exercise}[subsection]{Exercise}
\newtheorem{situation}[subsection]{Situation}

\theoremstyle{remark}
\newtheorem{remark}[subsection]{Remark}
\newtheorem{remarks}[subsection]{Remarks}

\numberwithin{equation}{subsection}

% Macros
%
\def\lim{\mathop{\mathrm{lim}}\nolimits}
\def\colim{\mathop{\mathrm{colim}}\nolimits}
\def\Spec{\mathop{\mathrm{Spec}}}
\def\Hom{\mathop{\mathrm{Hom}}\nolimits}
\def\Ext{\mathop{\mathrm{Ext}}\nolimits}
\def\SheafHom{\mathop{\mathcal{H}\!\mathit{om}}\nolimits}
\def\SheafExt{\mathop{\mathcal{E}\!\mathit{xt}}\nolimits}
\def\Sch{\mathit{Sch}}
\def\Mor{\mathop{\mathrm{Mor}}\nolimits}
\def\Ob{\mathop{\mathrm{Ob}}\nolimits}
\def\Sh{\mathop{\mathit{Sh}}\nolimits}
\def\NL{\mathop{N\!L}\nolimits}
\def\CH{\mathop{\mathrm{CH}}\nolimits}
\def\proetale{{pro\text{-}\acute{e}tale}}
\def\etale{{\acute{e}tale}}
\def\QCoh{\mathit{QCoh}}
\def\Ker{\mathop{\mathrm{Ker}}}
\def\Im{\mathop{\mathrm{Im}}}
\def\Coker{\mathop{\mathrm{Coker}}}
\def\Coim{\mathop{\mathrm{Coim}}}

% Boxtimes
%
\DeclareMathSymbol{\boxtimes}{\mathbin}{AMSa}{"02}

%
% Macros for moduli stacks/spaces
%
\def\QCohstack{\mathcal{QC}\!\mathit{oh}}
\def\Cohstack{\mathcal{C}\!\mathit{oh}}
\def\Spacesstack{\mathcal{S}\!\mathit{paces}}
\def\Quotfunctor{\mathrm{Quot}}
\def\Hilbfunctor{\mathrm{Hilb}}
\def\Curvesstack{\mathcal{C}\!\mathit{urves}}
\def\Polarizedstack{\mathcal{P}\!\mathit{olarized}}
\def\Complexesstack{\mathcal{C}\!\mathit{omplexes}}
% \Pic is the operator that assigns to X its picard group, usage \Pic(X)
% \Picardstack_{X/B} denotes the Picard stack of X over B
% \Picardfunctor_{X/B} denotes the Picard functor of X over B
\def\Pic{\mathop{\mathrm{Pic}}\nolimits}
\def\Picardstack{\mathcal{P}\!\mathit{ic}}
\def\Picardfunctor{\mathrm{Pic}}
\def\Deformationcategory{\mathcal{D}\!\mathit{ef}}
\usepackage{cleveref}
\crefname{originalreference}{original reference}{original references}
\Crefname{originalreference}{Original reference}{Original references}
\makeatletter
\newlabel{lemma-base-change-ind-etale}{{097J}{source}{}{sourcepointer-1}{}}
\newlabel{lemma-base-change-ind-etale@cref}{{[originalreference][0][]097J}{[1][0][]source}{}{}{}}
\makeatother
\crefname{originalreference}{original reference}{original references}
\Crefname{originalreference}{Original reference}{Original references}
\makeatletter
\newlabel{lemma-composition-ind-etale}{{097K}{source}{}{sourcepointer-2}{}}
\newlabel{lemma-composition-ind-etale@cref}{{[originalreference][0][]097K}{[1][0][]source}{}{}{}}
\makeatother
\crefname{originalreference}{original reference}{original references}
\Crefname{originalreference}{Original reference}{Original references}
\makeatletter
\newlabel{lemma-fully-faithful}{{0241}{source}{}{sourcepointer-3}{}}
\newlabel{lemma-fully-faithful@cref}{{[originalreference][0][]0241}{[1][0][]source}{}{}{}}
\makeatother
\crefname{originalreference}{original reference}{original references}
\Crefname{originalreference}{Original reference}{Original references}
\makeatletter
\newlabel{lemma-limit-pullback}{{099S}{source}{}{sourcepointer-4}{}}
\newlabel{lemma-limit-pullback@cref}{{[originalreference][0][]099S}{[1][0][]source}{}{}{}}
\makeatother
\crefname{originalreference}{original reference}{original references}
\Crefname{originalreference}{Original reference}{Original references}
\makeatletter
\newlabel{lemma-proetale-affine}{{098D}{source}{}{sourcepointer-5}{}}
\newlabel{lemma-proetale-affine@cref}{{[originalreference][0][]098D}{[1][0][]source}{}{}{}}
\makeatother
\crefname{originallemma}{original lemma}{original lemmas}
\Crefname{originallemma}{Original lemma}{Original lemmas}
\makeatletter
\newlabel{limits-lemma-descend-etale}{{07RP}{source}{}{sourcepointer-6}{}}
\newlabel{limits-lemma-descend-etale@cref}{{[originallemma][0][]07RP}{[1][0][]source}{}{}{}}
\makeatother
\crefname{originallemma}{original lemma}{original lemmas}
\Crefname{originallemma}{Original lemma}{Original lemmas}
\makeatletter
\newlabel{limits-lemma-descend-finite-presentation}{{01ZM}{source}{}{sourcepointer-7}{}}
\newlabel{limits-lemma-descend-finite-presentation@cref}{{[originallemma][0][]01ZM}{[1][0][]source}{}{}{}}
\makeatother
\crefname{originallemma}{original lemma}{original lemmas}
\Crefname{originallemma}{Original lemma}{Original lemmas}
\makeatletter
\newlabel{limits-lemma-descend-surjective}{{07RR}{source}{}{sourcepointer-8}{}}
\newlabel{limits-lemma-descend-surjective@cref}{{[originallemma][0][]07RR}{[1][0][]source}{}{}{}}
\makeatother
\crefname{originallemma}{original lemma}{original lemmas}
\Crefname{originallemma}{Original lemma}{Original lemmas}
\makeatletter
\newlabel{limits-lemma-limit-affine}{{01Z6}{source}{}{sourcepointer-9}{}}
\newlabel{limits-lemma-limit-affine@cref}{{[originallemma][0][]01Z6}{[1][0][]source}{}{}{}}
\makeatother
\crefname{originalreference}{original reference}{original references}
\Crefname{originalreference}{Original reference}{Original references}
\makeatletter
\newlabel{proposition-weakly-etale}{{097Z}{source}{}{sourcepointer-10}{}}
\newlabel{proposition-weakly-etale@cref}{{[originalreference][0][]097Z}{[1][0][]source}{}{}{}}
\makeatother
\crefname{originallemma}{original lemma}{original lemmas}
\Crefname{originallemma}{Original lemma}{Original lemmas}
\makeatletter
\newlabel{sites-cohomology-lemma-cech-vanish-collection}{{03F9}{source}{}{sourcepointer-11}{}}
\newlabel{sites-cohomology-lemma-cech-vanish-collection@cref}{{[originallemma][0][]03F9}{[1][0][]source}{}{}{}}
\makeatother
\crefname{originallemma}{original lemma}{original lemmas}
\Crefname{originallemma}{Original lemma}{Original lemmas}
\makeatletter
\newlabel{sites-cohomology-lemma-injective-trivial-cech}{{03AW}{source}{}{sourcepointer-12}{}}
\newlabel{sites-cohomology-lemma-injective-trivial-cech@cref}{{[originallemma][0][]03AW}{[1][0][]source}{}{}{}}
\makeatother
\title{Higher pro-etale cohomology of the pullback of an injective sheaf on an affine scheme vanishes}
\author{The Stacks Project Authors}
\date{stacks-project:proetale.tex:local-snapshot-20261012; GNU FDL 1.2 or later}
\begin{document}
\maketitle
\noindent\textit{Editorial note.} All delivered mathematical text is the author's own, copied from the pinned Stacks Project chapter file proetale.tex; no statement or proof step was rewritten, reordered or removed.\par\medskip
\noindent\textit{Editorial note.} The counted proof is the author's complete argument: the verification of the first two conditions of the Cech criterion, the reduction of a pro-etale covering to a filtered colimit of etale coverings, the approximation of an etale covering along an inverse limit of etale neighbourhoods, and the closing acyclicity argument.\par\medskip
\noindent\textit{Editorial note.} The label lemma-affine-vanishing names the counted statement, and the proof environment immediately following it in the source is delivered verbatim as its proof.\par\medskip
\noindent\textit{Editorial note.} The counted statement is carried by the author's own lemma environment in the chapter source.\par\medskip
\noindent\textit{Editorial note.} Pre-open gate: python operations/agent-authoring/tools/precard\_gate.py --source proetale.tex --label lemma-affine-vanishing --start 4082 --end 4167 --local-id 3576 returned RESULT=OK (L3576 free, occupancy FREE, public interval clear). Candidate taken from sources/stacks-proetale/candidates-halfA.json (new exclusive pool assigned by the Lead on 2026-10-13).\par\medskip
\noindent\textit{Editorial note.} The statement and proof are copied from the pinned chapter file; the cited inputs are kept as external pointers with their canonical Stacks tags.\par\medskip
\noindent\textit{Allowed standard prerequisites (editorial).}\begin{itemize}
\item The Cech vanishing criterion for a collection of objects and coverings. Cohomology on Sites, the lemma giving conditions (1), (2), (3) on a set of objects and a set of coverings under which the higher Cech cohomology of an injective sheaf vanishes. Supplies the criterion applied to the set of affine objects whose structure rings are ind-etale over the ring of X and to the pro-etale coverings between them.
\item Permanence and base change for ind-etale algebras, and affineness of the pro-etale site. Same chapter, the lemmas on composition and base change of ind-etale algebras and the lemma that an affine scheme is a final object with the expected coverings in the pro-etale topology, together with the proposition on weakly etale morphisms. Supply conditions (1) and (2) of the criterion.
\item Comparison of the pullback along the pro-etale to etale morphism with restriction. Same chapter, the lemma identifying the pullback of an etale sheaf to the pro-etale site with its restriction, and the lemma computing finite limits of the resulting sections. Supply the identification of the Cech complexes of the pullback with those of the original sheaf over the approximating etale coverings.
\item Acyclicity of the Cech complex of an injective sheaf. Cohomology on Sites, the lemma that the Cech complex of an injective sheaf over a covering is acyclic in positive degrees. Supplies the vanishing of the Cech cohomology of the approximating etale coverings, which passes to the limit.
\item Descent of finite presentation, affine schemes and etale morphisms along a directed limit. Limits, the lemmas descending finite presentation, affine schemes and etale morphisms along a directed inverse limit. Supply the approximation of the etale covering of U by an etale covering of one U\_b.
\end{itemize}\medskip
\ifdefined\hypertarget\hypertarget{sourcepointer-1}{}\fi
\noindent\textit{Source reference.} Original reference 097J: sources/stacks-proetale/pinned-original-source.tex, source lines 1387--1387. Prerequisite pointer invoked at proof line(s) 4101; it names a standard fact of the same chapter that the counted argument uses, and every step of the counted proof built on it is retained verbatim in the delivered body.\par\medskip
\ifdefined\hypertarget\hypertarget{sourcepointer-2}{}\fi
\noindent\textit{Source reference.} Original reference 097K: sources/stacks-proetale/pinned-original-source.tex, source lines 1398--1398. Prerequisite pointer invoked at proof line(s) 4100; it names a standard fact of the same chapter that the counted argument uses, and every step of the counted proof built on it is retained verbatim in the delivered body.\par\medskip
\ifdefined\hypertarget\hypertarget{sourcepointer-3}{}\fi
\noindent\textit{Source reference.} Original reference 0241: sources/stacks-proetale/pinned-original-source.tex, source lines 4045--4045. Prerequisite pointer invoked at proof line(s) 4159; it names a standard fact of the same chapter that the counted argument uses, and every step of the counted proof built on it is retained verbatim in the delivered body.\par\medskip
\ifdefined\hypertarget\hypertarget{sourcepointer-4}{}\fi
\noindent\textit{Source reference.} Original reference 099S: sources/stacks-proetale/pinned-original-source.tex, source lines 4058--4058. Prerequisite pointer invoked at proof line(s) 4129; it names a standard fact of the same chapter that the counted argument uses, and every step of the counted proof built on it is retained verbatim in the delivered body.\par\medskip
\ifdefined\hypertarget\hypertarget{sourcepointer-5}{}\fi
\noindent\textit{Source reference.} Original reference 098D: sources/stacks-proetale/pinned-original-source.tex, source lines 2433--2433. Prerequisite pointer invoked at proof line(s) 4102; it names a standard fact of the same chapter that the counted argument uses, and every step of the counted proof built on it is retained verbatim in the delivered body.\par\medskip
\ifdefined\hypertarget\hypertarget{sourcepointer-6}{}\fi
\noindent\textit{Source reference.} Original lemma 07RP: sources/stacks-chow/tags.txt, source lines 9874--9874. Prerequisite pointer invoked at proof line(s) 4144; it names a standard fact of a sibling chapter that the counted argument uses, and every step of the counted proof built on it is retained verbatim in the delivered body.\par\medskip
\ifdefined\hypertarget\hypertarget{sourcepointer-7}{}\fi
\noindent\textit{Source reference.} Original lemma 01ZM: sources/stacks-chow/tags.txt, source lines 2599--2599. Prerequisite pointer invoked at proof line(s) 4142; it names a standard fact of a sibling chapter that the counted argument uses, and every step of the counted proof built on it is retained verbatim in the delivered body.\par\medskip
\ifdefined\hypertarget\hypertarget{sourcepointer-8}{}\fi
\noindent\textit{Source reference.} Original lemma 07RR: sources/stacks-chow/tags.txt, source lines 9876--9876. Prerequisite pointer invoked at proof line(s) 4152; it names a standard fact of a sibling chapter that the counted argument uses, and every step of the counted proof built on it is retained verbatim in the delivered body.\par\medskip
\ifdefined\hypertarget\hypertarget{sourcepointer-9}{}\fi
\noindent\textit{Source reference.} Original lemma 01Z6: sources/stacks-chow/tags.txt, source lines 2581--2581. Prerequisite pointer invoked at proof line(s) 4143; it names a standard fact of a sibling chapter that the counted argument uses, and every step of the counted proof built on it is retained verbatim in the delivered body.\par\medskip
\ifdefined\hypertarget\hypertarget{sourcepointer-10}{}\fi
\noindent\textit{Source reference.} Original reference 097Z: sources/stacks-proetale/pinned-original-source.tex, source lines 1802--1802. Prerequisite pointer invoked at proof line(s) 4103; it names a standard fact of the same chapter that the counted argument uses, and every step of the counted proof built on it is retained verbatim in the delivered body.\par\medskip
\ifdefined\hypertarget\hypertarget{sourcepointer-11}{}\fi
\noindent\textit{Source reference.} Original lemma 03F9: sources/stacks-chow/tags.txt, source lines 4484--4484. Prerequisite pointer invoked at proof line(s) 4091; it names a standard fact of a sibling chapter that the counted argument uses, and every step of the counted proof built on it is retained verbatim in the delivered body.\par\medskip
\ifdefined\hypertarget\hypertarget{sourcepointer-12}{}\fi
\noindent\textit{Source reference.} Original lemma 03AW: sources/stacks-chow/tags.txt, source lines 4325--4325. Prerequisite pointer invoked at proof line(s) 4163; it names a standard fact of a sibling chapter that the counted argument uses, and every step of the counted proof built on it is retained verbatim in the delivered body.\par\medskip
\begin{lemma}
\label{lemma-affine-vanishing}
Let $X$ be an affine scheme. For injective abelian sheaf $\mathcal{I}$ on
$X_\etale$ we have $H^p(X_\proetale, \epsilon^{-1}\mathcal{I}) = 0$
for $p > 0$.
\end{lemma}

\begin{proof}
We are going to use
Cohomology on Sites, Lemma \ref{sites-cohomology-lemma-cech-vanish-collection}
to prove this. Let $\mathcal{B} \subset \Ob(X_\proetale)$ be the set of affine
objects $U$ of $X_\proetale$ such that $\mathcal{O}(X) \to \mathcal{O}(U)$ is
ind-\'etale. Let $\text{Cov}$ be the set of pro-\'etale coverings
$\{U_i \to U\}_{i = 1, \ldots, n}$ with $U \in \mathcal{B}$ such that
$\mathcal{O}(U) \to \mathcal{O}(U_i)$ is ind-\'etale for $i = 1, \ldots, n$.
Properties (1) and (2) of
Cohomology on Sites, Lemma \ref{sites-cohomology-lemma-cech-vanish-collection}
hold for $\mathcal{B}$ and $\text{Cov}$ by
Lemmas \ref{lemma-composition-ind-etale},
\ref{lemma-base-change-ind-etale}, and
\ref{lemma-proetale-affine} and
Proposition \ref{proposition-weakly-etale}.

\medskip\noindent
To check condition (3) suppose that
$\mathcal{U} = \{U_i \to U\}_{i = 1, \ldots, n}$
is an element of $\text{Cov}$. We have to show that the higher
Cech cohomology groups of $\epsilon^{-1}\mathcal{I}$
with respect to $\mathcal{U}$ are zero.
First we write $U_i = \lim_{a \in A_i} U_{i, a}$ as a directed inverse limit
with $U_{i, a} \to U$ \'etale and $U_{i, a}$ affine.
We think of $A_1 \times \ldots \times A_n$ as a direct set
with ordering $(a_1, \ldots, a_n) \geq (a_1', \ldots, a_n')$
if and only if $a_i \geq a_i'$ for $i = 1, \ldots, n$.
Observe that
$\mathcal{U}_{(a_1, \ldots, a_n)} = \{U_{i, a_i} \to U\}_{i = 1, \ldots, n}$
is an \'etale covering for all
$a_1, \ldots, a_n \in A_1 \times \ldots \times A_n$.
Observe that
$$
U_{i_0} \times_U U_{i_1} \times_U \ldots \times_U U_{i_p}
=
\lim_{(a_1, \ldots, a_n) \in A_1 \times \ldots \times A_n}
U_{i_0, a_{i_0}} \times_U
U_{i_1, a_{i_1}} \times_U \ldots \times_U U_{i_p, a_{i_p}}
$$
for all $i_0, \ldots, i_p \in \{1, \ldots, n\}$ because limits
commute with fibred products. Hence by Lemma \ref{lemma-limit-pullback}
and exactness of filtered colimits we have
$$
\check{H}^p(\mathcal{U}, \epsilon^{-1}\mathcal{I}) =
\colim \check{H}^p(\mathcal{U}_{(a_1, \ldots, a_n)}, \epsilon^{-1}\mathcal{I}) 
$$
Thus it suffices to prove the vanishing for \'etale coverings
of $U$!

\medskip\noindent
Let $\mathcal{U} = \{U_i \to U\}_{i = 1, \ldots, n}$ be an \'etale covering
with $U_i$ affine. Write $U = \lim_{b \in B} U_b$ as a directed inverse limit
with $U_b$ affine and $U_b \to X$ \'etale.
By Limits, Lemmas \ref{limits-lemma-descend-finite-presentation},
\ref{limits-lemma-limit-affine}, and
\ref{limits-lemma-descend-etale}
we can choose a $b_0 \in B$ such that for $i = 1, \ldots, n$ there is an
\'etale morphism $U_{i, b_0} \to U_{b_0}$ of affines
such that $U_i = U \times_{U_{b_0}} U_{i, b_0}$.
Set $U_{i, b} = U_b \times_{U_{b_0}} U_{i, b_0}$ for $b \geq b_0$.
For $b$ large enough the family
$\mathcal{U}_b = \{U_{i, b} \to U_b\}_{i = 1, \ldots, n}$
is an \'etale covering, see
Limits, Lemma \ref{limits-lemma-descend-surjective}.
Exactly as before we find that
$$
\check{H}^p(\mathcal{U}, \epsilon^{-1}\mathcal{I}) =
\colim \check{H}^p(\mathcal{U}_b, \epsilon^{-1}\mathcal{I}) =
\colim \check{H}^p(\mathcal{U}_b, \mathcal{I})
$$
the final equality by Lemma \ref{lemma-fully-faithful}.
Since each of the {\v C}ech complexes on the right hand side is
acyclic in positive degrees
(Cohomology on Sites, Lemma
\ref{sites-cohomology-lemma-injective-trivial-cech})
it follows that the one on the left is too. This proves condition (3) of
Cohomology on Sites, Lemma \ref{sites-cohomology-lemma-cech-vanish-collection}.
Since $X \in \mathcal{B}$ the lemma follows.
\end{proof}

\end{document}
